\documentclass[11pt,a4paper]{article}
% Estilo comum dos PDFs individuais da Sprint 1.
\usepackage[utf8]{inputenc}
\usepackage[T1]{fontenc}
\usepackage[brazilian]{babel}
\usepackage[a4paper,margin=2.35cm]{geometry}
\usepackage{xcolor}
\usepackage{hyperref}
\usepackage{enumitem}
\usepackage{listings}
\usepackage{parskip}
\usepackage{needspace}
\definecolor{azul}{RGB}{20,63,112}
\hypersetup{colorlinks=true,linkcolor=azul,urlcolor=azul}
\setlist[itemize]{leftmargin=*,itemsep=3pt}
\lstset{basicstyle=\ttfamily\small,frame=single,breaklines=true,columns=fullflexible}
\sloppy
\newcommand{\codigo}[1]{\texttt{#1}}
\newcommand{\titulo}[2]{\begin{center}{\LARGE\bfseries #1\par}\vspace{4pt}{\large #2\par}\vspace{8pt}\rule{0.8\textwidth}{0.4pt}\par\vspace{5pt}{\small Matheus Moreli Durão Silva -- Matrícula 12621CBS217\par Cibersegurança - Bacharelado - Noturno\par}\end{center}}
\newcommand{\cardfonte}[1]{\par\smallskip\noindent\textit{Referência: }\href{#1}{Trello - enunciado e checklist do card}.}

\begin{document}\titulo{S1\_19 - Pull Request}{Revisar antes de integrar}
\section*{Entregável correto}
O card pede uma Pull Request completa, documentada e aprovada. PR não é apenas um botão de merge: é uma proposta de mudança, com título, descrição, diferenças legíveis e registro de revisão antes da integração.
\section*{Repositório remoto, clone e fork}
O repositório remoto vive no GitHub. \codigo{git clone} cria uma cópia local ligada ao remoto; \textit{fork} cria uma cópia do repositório na conta de outra pessoa, útil quando não há permissão direta para gravar no original. Após mudar arquivos localmente, commits e \codigo{push} disponibilizam uma branch no remoto para abrir a PR.
\section*{Como revisar uma PR}
Primeiro, conferir o objetivo: que problema a mudança resolve? Depois, ler a aba de arquivos alterados e procurar efeitos indesejados, como código fora do escopo ou arquivo de dados exposto. Comentários devem apontar uma linha e explicar a razão. Se há conflito com a base, ele deve ser resolvido antes ou durante o merge, seguindo a versão que preserva as duas intenções.
\section*{Aplicação no projeto}
As PRs de validação JSON e de resumo de inventário tornam visível a evolução do \codigo{app.py}. Essa documentação liga o que foi aprendido em Git a uma decisão concreta: uma mudança foi desenvolvida, comparada com \codigo{main}, revisada e integrada.
\section*{Checklist demonstrado}
Repositório remoto, clone/fork, criação de PR, revisão, integração e conflitos. Uma PR completa tem origem e destino claros, descrição objetiva, alterações verificáveis e resultado do merge registrado.
\cardfonte{https://trello.com/c/tsLFX3pi}
\end{document}
